% Folder 93, batch 08 — archivists' pages 141–160.
% Pass: Opus 5.5 (claude-opus-5-5), 2026-10-10 — first pass, unchecked against the pages by a human.
%
% Manuscript notes in Grothendieck's hand throughout. Connections on affine
% torsors, de Rham cohomology with coefficients, interpretation of H^1_DR by
% torsors with integrable connection; then stratifications.
%
\documentclass[11pt,a4paper]{article}
\input{../preamble/grothendieck}

\folder{93}
\batch{8}
\pages{141}{160}
\foldertitle{Connexions de Gauss-Manin et équations de Picard-Fuchs. Opération de Cartier : copies de tapuscrit annoté (s.d.), tiré à part (1981), notes manuscrites (s.d.).}
\dating{[1972]-1981}
\watermark{Édition de démonstration}

\begin{document}

\section{1) Connexions sur torseurs affines}

\page{141}
\note{le titre, souligné, est de sa main ; « 1) » ouvre une suite numérotée qui commence sur cette page.}

$S/T$ schéma relatif \add{($q : S \to T$)}, $E$ module quasi-cohérent sur $S$, d'où :
\[
(1)\qquad 0 \longrightarrow E \otimes \Omega^1_{S/T} \longrightarrow P^1_{S/T}(E) \longrightarrow E \longrightarrow 0
\]
extension de $\mathcal{O}_S$-modules. On a un splitting $q^{-1}(\mathcal{O}_T)$-linéaire
(mais pas $\mathcal{O}_S$-linéaire)
\[
(2)\qquad d^1_E : E \longrightarrow P^1_{S/T}(E)
\]
qui permet d'identifier $P^1_{S/T}(E)$ \emph{additivement} à
$E \oplus E \otimes \Omega^1_{S/T}$, en écrivant donc \struck{\ill{}} \add{les sections}
de $P^1_{S/T}(E)$ \add{(sur \ill{})} \uncertain{de façon unique} comme
\[
(3)\qquad \xi = d^1_E x + \omega \qquad
\begin{cases} \omega \in \Gamma(\ ,E \otimes \Omega^1_{S/T}) \\ x \in \Gamma(\ ,E) \end{cases}
\]
\note{sous le premier membre, une formule biffée et illisible, terminée par un indice ${}_{S/T}$.}

Comme, si $f$ est \uncertain{une section de} $\mathcal{O}_S$,
\[
f\xi = f\omega + f\,d^{(1)}_E x = f\omega + d^{(1)}_E(fx) + [f - d^{(1)}(f)]\,d^1_E(x)
\]
\note{sous $d^{(1)}_E(fx)$, il écrit (avec un signe d'égalité vertical) $d^1(f)\,d^1_E(x)$ ; sous le crochet, « $-df$ ».}
i.e.

\struck{(4) $f(\omega + d^{(1)}_E x) = f\omega$}
\[
(4)\qquad f(d^1_E x + \omega) = d^1_E fx + (f\omega - x \otimes df) .
\]
\note{le signe entre $f\omega$ et $x \otimes df$ est surchargé : un « $-$ » repassé, que la p.~142 confirme.}

Une \emph{connexion} $\delta_E$ de $E$ rel. à $T$ est un splitting \add{$\mathcal{O}_S$-linéaire}
de (1) ; \uncertain{donc tout splitting additif}, comme dans (2), il est donné par un
homomorphisme additif
\[
(5)\qquad d^0_E : E \longrightarrow E \otimes \Omega^1_{S/T}
\]
\ill{}
\[
(6)\qquad \delta_E x = d^1_E(x) - d^0_E(x) ,
\]
il faut expliciter la condition sur $d^0_E$ de (5) pour que $\delta_E$ \uncertain{soit}
$\mathcal{O}_S$-linéaire,

\page{142}
i.e. qu'on ait $\delta f x = f \delta x$, i.e. compte tenu de (4),
\[
d^1_E(fx) + (-f d^0_E x - x \otimes df) = d^1_E(fx) - d^0_E(fx) ,
\]
ce qui s'écrit
\[
(7)\qquad d^0_E(fx) = f d^0_E(x) + x \otimes df .
\]
On a aussi, utilisant la structure de $\Omega^\bullet_{S/T}$-module à droite de
$E \otimes \Omega^\bullet_{S/T}$ :
\[
(7')\qquad d^0_E(xf) = d^0_E(x) f + x\,df .
\]
\note{dans (7'), $xf$ est écrit sur un $fx$ biffé.}

NB On dit souvent que $d^0_E$ (5) satisfaisant (7) est une \emph{connexion} sur $E$ rel. à $T$.
\struck{On sait alors \ill{}}

On sait alors qu'il $\exists$ prolongement unique de $d^0_E$
\struck{en une dérivation \ill{}} en un endomorphisme additif
\[
(8)\qquad d^\bullet_E : E \otimes \Omega^\bullet_{S/T} \longrightarrow E \otimes \Omega^\bullet_{S/T}
\]
qui \struck{\ill{}} satisfasse la relation suivante rel. à la structure de
$\Omega^\bullet_{S/T}$-module à droite :
\[
(9)\qquad d^\bullet_E(\varphi\omega) = d_E(\varphi)\,\omega + (-1)^{\deg\varphi}\,\varphi\, d(\omega) ,
\]
et où $d^\bullet_E$ est un \uncertain{opérateur} de degré 1. Pour que l'on ait
\[
(10)\qquad d^\bullet_E \circ d^\bullet_E = 0 ,
\]
il faut et il suffit que l'on ait

(10') le composé $d^1_E d^0_E : E \to E \otimes \Omega^1_{S/T} \to E \otimes \Omega^2_{S/T}$
est nul.

Le composé en question, considéré comme section de
$\mathcal{H}om(E, E \otimes \Omega^2_{S/T})$, est appelé la \emph{courbure} de la connexion.
La connexion est dite \emph{intégrable} si la courbure est nulle, i.e. si $d^\bullet_E$ (8)
fait

\page{143}
de $E \otimes \Omega^\bullet_{S/T}$ un complexe. Supposons qu'il en soit ainsi.
\struck{\ill{}}
\[
(11)\qquad H^*_{DR}(S,E) = \mathbb{H}^*(S, E \otimes \Omega^\bullet_{S/T}) .
\]
\note{sous $E$, une flèche renvoie à : « sous-entendu : $E$ muni de $\delta_E$, i.e. de $d^0_E$ ».}

On a une suite spectrale
\[
(12)\qquad H^*_{DR}(S,E) \Longleftarrow E_1^{pq} = H^q(S, E \otimes \Omega^p_{S/T})
\]
\note{devant $H^q$, une première écriture biffée et illisible.}
associée à la filtration triviale du complexe $E \otimes \Omega^\bullet_{S/T}$.

On a aussi des homomorphismes \ill{}
\[
(13)\qquad H^*(\Gamma(S, E \otimes \Omega^\bullet_{S/T})) \Longrightarrow H^*_{DR}(S,E)
\]
qui sont des isom. si p.ex. $S$ affine, et
\[
(14)\qquad H^i_{DR}(S,E) \longrightarrow H^i(S,E) ,
\]
qui \ill{} \uncertain{provient} des homomorphismes de complexes \uncertain{évidents} (troncation) :
\[
(15)\qquad E \otimes \Omega^\bullet_{S/T} \longrightarrow E .
\]

NB Pour $E$ \ill{} : conn. intégrable \uncertain{variable},
$H^*_{DR}(S,E)$ est un $\partial$-foncteur en $E$ (cat. \ill{} aussi : cat. abélienne). Les
suites spectrales \ill{} et les hom. écrits plus haut sont fonctoriels en $E$.

Re NB \struck{I} Images inverses \add{(pour un diagramme commutatif}
\begin{tikzcd}[column sep=small, row sep=small]
S \arrow[d] & S' \arrow[l] \arrow[d] \\
T & T' \arrow[l]
\end{tikzcd}
\add{\ill{})} tout suit \ill{} !

En particulier : \struck{\ill{}} a) Image inverse par $T' \to T$ (on prend $S' = S \times_T T'$).
b) Si $S \to T' \to T$, une connexion rel. à $T$ en définit une rel. à $T'$, et

\page{144}
\note{la phrase de la p.~143 reste en suspens ; la page reprend sur un nouvel alinéa.}

Nous allons étudier (14) pour $i = 0, 1$.

1) $H^0_{DR}(S,E)$ \uncertain{est formé} des sections de $E$ qui sont
\struck{\ill{}} « horizontales pour la connexion $\delta$ », ce qui \uncertain{signifie}
\struck{$d^0_E \ldots = 0$}
\[
(16)\qquad d^0_E(x) = 0 \qquad\text{i.e.}\qquad \delta_E x = d^1_E x
\]
Alors (14) est l'injection canonique. \struck{\ill{}}

\marginal{\ill{} cf. \uncertain{additions} 1.8.}
2) Soit $K^\bullet : 0 \to K^0 \xrightarrow{d^0} K^1 \xrightarrow{d^1} K^2 \to \cdots$ un complexe
\struck{\ill{}} de faisceaux abéliens sur un \uncertain{topos} $S$.
\struck{On constate que $\mathbb{H}^1 S$} On constate que $\mathbb{H}^1(S,K^\bullet)$
s'interprète comme le groupe des classes d'isomorphie d'objets de la catégorie suivante
$C(K^\bullet)$ :

$\mathrm{Ob}\,C(K^\bullet) =$ couples $(P^\bullet, \xi)$, avec $P$ un torseur sous $K^0$,
$\xi$ une section de $P^1 = P \times^{K^0} K^1$, telle que la section
\struck{\ill{}} $d^1_*(\xi)$ de $P^2 = P \times^{K^0} K^2 \simeq P^1 \times^{K^1} K^2$
soit la section canonique, \uncertain{déduite} du fait que $d^1 d^0 = 0$.

$\mathrm{Fl}\,C(K^\bullet)$ : on devine.

\struck{Cette description de $\mathbb{H}^1(S,K) \to \mathbb{H}^1(S\ldots$}
\struck{L'homom. canonique} Cette description de $\mathbb{H}^1(S,K^\bullet)$ \ill{}
en \uncertain{tous} les complexes $K^\bullet$. Par suite,
l'hom. \add{can.} \struck{$\mathbb{H}(S,K^\bullet$}
\[
\mathbb{H}^1(S,K^\bullet) \longrightarrow H^1(S,K^0)
\]
déduit de l'\struck{\ill{}} troncation $K^\bullet \to K^0$ correspond au foncteur
$(P^0,\xi) \longmapsto P^0$ « oubli de $\xi$ ».

Notons que \struck{\ill{}} \add{\ill{}} $P^0$ donné (torseur sous $K^0$), la donnée d'une
trivialisation $\xi$ de $P^1$ \uncertain{équivaut} à un isom. \struck{\ill{}} \add{\ill{}}
$P^1 \simeq K^1$, donc les $\xi$ correspondent

\page{145}
aux homom.
\[
\bar d^0 : P^0 \longrightarrow K^1
\]
satisfaisant
\[
\bar d^0(xf) = \bar d^0(x)\, d^0(f) \qquad x \in P^0(\cdot),\ f \in K^0(\cdot) .
\]
\note{ainsi sur la page, sans signe entre $\bar d^0(x)$ et $d^0(f)$ : la loi de groupe de $K^1$ est notée multiplicativement.}
La condition mise ci-dessus sur $\xi$ signifie simplement que $\bar d^0$ prend ses valeurs
dans $Z^1K = \mathrm{Ker}(d^1)$.

Ceci s'applique \uncertain{tel quel} à la description de $H^1_{DR}(S,E)$, comme
\uncertain{classifiant} les objets suivants : Un torseur $P$ sous $E$, un hom. de faisceaux
d'ensembles
\[
(17)\qquad d_P : P \longrightarrow (E \otimes \Omega^1_{S/T}) ,
\]
\note{devant $(E \otimes \Omega^1_{S/T})$, un $Z$ biffé par hachures.}
satisfaisant aux conditions
\[
(18)\qquad
\begin{cases}
d_P(xf) = d_P(x)\, d^0_E(f) & x \in P(\cdot),\ f \in \mathcal{O}_S(\cdot) . \\
d_P(P) \subset Z^1(E \otimes \Omega_{S/T}) & \text{i.e. } d^1_E d_P = 0 .
\end{cases}
\]
\note{« $f \in \mathcal{O}_S(\cdot)$ » est lu tel quel ; on attendrait $f \in E(\cdot)$.}
Un tel $d_P$ s'appellera une \emph{connexion} \add{\emph{intégrable}} sur le torseur $P$,
\emph{compatible avec la connexion} \add{\emph{intégrable}} $d^0_E$ sur $E$.

\emph{NB} Si $d^0_E$ est une connexion sur $E$ (pas nécess. intégrable),
\struck{\ill{}} on peut se \uncertain{proposer} de \uncertain{trouver}, pour un torseur $P$ sous $E$,
quelles sont les connexions compatibles avec $d^0_E$, et \uncertain{plus}
généralement \struck{\ill{}} de \ill{} \uncertain{voisinages infinitésimaux}
d'indice 1. Je \uncertain{dis} que si $S$ est affine, il existe
de \uncertain{telles} connexions, donc il en existe toujours localement,
et que les connexions forment un torseur sous
$E \otimes \Omega$, qui en fait est isomorphe à $P \times^E (E \otimes \Omega)$
(grâce à $d^0_E : E \to E \otimes \Omega$). Donc les connexions

\page{146}
On trouve en particulier une interprétation de $H^1_{DR}(S,E)$, comme classifiant les
torseurs $P$ sous $E$, munis d'une
\note{la page s'arrête là ; le reste du feuillet est blanc. La phrase de la p.~145 (« Donc les connexions ») se poursuit en tête de la p.~147.}

\page{147}
s'interprètent comme les hom (17) satisfaisant la \emph{première relation} (18).
\struck{C'est \ill{}} C'est la deuxième relation qui précise que l'on se restreint aux
connexions \emph{intégrables} sur $P$.

Les \struck{faisceaux des} automorphismes de $P$ \emph{compatibles avec} $d_P$
\struck{\ill{}} et \struck{\ill{}} sont \uncertain{donnés} par \struck{$x \in \Gamma(S,E$}
\ill{} du gpe $\Gamma(S,E)$ des \uncertain{automorphismes} \uncertain{fournis} par les sections
\emph{horizontales} de $E$, i.e. $H^0_{DR}(S,E)$.

Les descriptions qui précèdent \add{fonctorielles en $E$, et} \add{compatibles}
aux images inverses \uncertain{envisagées} (cf.) plus haut.

\section{2) Stratifications sur torseurs affines}

Supposons donné sur $E$ une \uncertain{stratification} rel. à $T$. \add{d'où un faisceau
$E_{\mathrm{strat}}$} sur $(S/T)_{\mathrm{strat}}$, donc \ill{} \uncertain{par} \ill{}
\[
(19)\qquad H^*_{\mathrm{strat}}(S,E) \overset{\mathrm{déf}}{=} H^*((S/T)_{\mathrm{strat}}, E_{\mathrm{strat}}) .
\]
\struck{La stratification de $E$} \struck{On a défini des $[\ ]$ \ill{} naturel}

Une donnée de \uncertain{voisinages} différentiels d'ordre \uncertain{$\infty$} sur $E$ revient
à la donnée d'un homomorphisme
\[
(20)\qquad \delta^\infty_E : E \longrightarrow P^\infty_{S/T}(E)
\]
\marginal{On dit parfois \uncertain{que} « $\delta^\infty_E$ » \uncertain{connexion}.}
satisfaisant la condition \uncertain{que} $P^\infty_{S/T}(E) \xrightarrow{\varepsilon_E} E$ \uncertain{donne} l'identité :
$\varepsilon_E \delta^\infty_E : E \to P^\infty_{S/T}(E) \to E$.

(21) On peut \ill{} \struck{\ill{}}
$d^+_E : E \to P^\infty_{S/T}(E)^+$, $d^+_E(x) = d^\infty(x) - \delta^\infty_E(x)$,
\struck{\ill{}}

\struck{(21) La condition \ill{} stratification \ill{} une condition sur $\delta_E$ qui}
\marginal{$P^\infty_{S/T}(E)^+$, idéal d'augmentation.}
\marginal{(plus bas dans no 3.)}

\page{148}
\struck{implique que} le composé de (21) avec
\struck{$P^\infty_{S/T}$}
\[
P^\infty_{S/T}(E)^+ \longrightarrow P^1_{S/T}(E)^+ \simeq E \otimes \Omega^1_{S/T} ,
\]
\struck{\ill{}} noté $d^0_E$
\[
(23)\qquad d^0_E : E \longrightarrow E \otimes \Omega^1_{S/T} , \qquad d^0_E = d^+_E \bmod \ldots
\]
\note{le numéro lu « (23) » pourrait être « (22) » ; sous le $d^0_E$ de droite, un signe surchargé.}
est l'hom. (5) associé à la connexion de $E$ définie par la donnée de voisinages
\uncertain{infinitésimaux}. \struck{Si} \emph{On vérifie que si} \struck{infinitésimale \ill{}} $\delta_E$ est une
stratification, \uncertain{alors} la donnée de voisinages \struck{infinitésimaux}
dans la connexion associée est intégrable, i.e. (23)
définit un complexe de De Rham.
On sait bien [ ] le \ill{} \ill{} des hom. canoniques
\[
(24)\qquad H^*_{\mathrm{strat}}(S,E) \longrightarrow H^*_{DR}(S,E) ,
\]
\uncertain{compatibles} \uncertain{aux changements de base}.
\note{les deux lignes « compatibles … base » sont reliées par un crochet à (24).}

\struck{\uncertain{Intervenant} en $F$,} Nous voulons une interprétation de $H^1_{\mathrm{strat}}(S,E)$
analogue : \uncertain{celle} \uncertain{donnée} \uncertain{plus haut} pour $H^1_{DR}(S,E)$.

\note{la suite de la page est biffée par des hachures obliques et une grande croix ; on la transcrit comme biffée.}
\struck{Pour ceci, \ill{} Soit, \ill{} $C^\bullet(E)$ \ill{} \ill{} de $E$ \ill{}
\ill{} le faisceau sur $S$ \ill{} \ill{} \ill{} le $X$ \ill{} de la \ill{} de $(S/T)^{\ill{}}$ de l'\ill{}
formel \ill{} $E$ \ill{} \ill{} (\ill{}) \ill{}. La stratification de}
\struck{(24) $C^0(E) = E$, $C^1(E) = P^\infty_{S/T}(E)$.}
\struck{La stratification de $E$ met sur $C^\bullet(E)$ une structure de faisceau
cosimplicial \ill{} [en \uncertain{faisceaux} en comodules sur $C^\bullet(\mathcal{O}_S)$] :}
\note{la suite biffée porte encore : (25) $C^0(E) \rightrightarrows C^1(E) \Rrightarrow C^2(E) \ldots$, avec dessous $C^0(E) = E \rightrightarrows P^\infty_{S/T}(E)$ par les deux flèches $\delta^\infty_E$ et $d^\infty_E$ ; à gauche, un essai de flèche $E \to P^\infty_{S/T}(E)$, raturé.}

\page{149}
\marginal{$U \subset U'$ au-dessus de $S \to T$, petit diagramme.}
Or un torseur sous $S_{\mathrm{strat}}$ \add{sur} $E_{\mathrm{strat}}$ peut être interprété
comme la donnée, pour tout objet \add{$(U,U')$ de} $S_{\mathrm{strat}}$, d'un
torseur sous $(E_{\mathrm{strat}})_{(U,U')} = E_{U'}$, \struck{\ill{}} plus \ill{} de façon
contravariante en $(U,U')$. On voit que ceci équivaut
à la donnée d'un torseur ordinaire $P$ sous $E$, et
d'une \emph{stratification} sur $P$, compatible avec celle de $E$.

\struck{Nous regardons} Nous regardons $P^\infty_{S/T}(E)$ comme un \emph{pro-faisceau}
\add{A.R.} sur $S$, et comme tel il est muni d'une stratification
canonique (\uncertain{indépendante} de celle de $E$, cf.~[ ]).
Nous \uncertain{avons} là \struck{\ill{}}, suite exacte
\[
(25)\qquad 0 \longrightarrow E \xrightarrow{\ \delta = \delta^\infty_E\ } P^\infty_{S/T}(E) \xrightarrow{\ p\ } P^\infty_{S/T}(E)/\delta(E) \longrightarrow 0 ,
\]
où le dernier terme est défini comme le \uncertain{support}
\uncertain{rigide} des \struck{\ill{}} \add{\uncertain{conoyaux}} des \uncertain{composés}
$E \xrightarrow{\delta_E} P^\infty_{S/T}(E) \to P^n_{S/T}(E)$.
Notons \struck{\ill{}} que l'hom. $\delta$ est compatible avec les
stratifications. La donnée d'un torseur stratifié $P_{\mathrm{strat}}$
sous $E$ s'identifie donc : à la donnée d'un
torseur stratifié sous le pro-faisceau $P^\infty_{S/T}(E)$, et
d'une trivialisation \struck{\ill{}} du \struck{\ill{}} torseur stratifié associé
du groupe $P^\infty_{S/T}(E)$. On constate qu'un torseur
stratifié sous $P^\infty_{S/T}(E)$, c'est « la même chose » qu'un
torseur ordinaire sous $E$ (cf.~[ ] 6.4), à un
tel $P$ étant associé le pro-faisceau \add{$P^\infty_{S/T}(P)$} \struck{des \ill{}}
\uncertain{jets} de $P$ sur $S/T$ d'ordre $\infty$ (\ill{}). Alors
la trivialisation $P^\infty_{S/T}(P) \to P^\infty_{S/T}(E)/\delta E$ revient à la

\page{150}
donnée d'un hom. de faisceaux
\[
(26)\qquad P \xrightarrow{\ \lambda\ } P^\infty_{S/T}(E)/\delta E
\]
\note{l'indice de $\lambda$ en (26) est raturé.}
satisfaisant la condition
\[
\text{\struck{$\lambda(x + \xi) = \lambda(x) + d^\infty_{S/T}(\xi)$}}
\]
\[
(27)\qquad \lambda(x + \xi) = \lambda(x) + p(d^\infty_{S/T}(\xi)) \qquad x \in P(\cdot),\ \xi \in E(\cdot) .
\]
Il faut encore exprimer que la trivialisation du
$P^\infty_{S/T}(E)/E$-torseur obtenu est compatible avec les
stratifications. Comme $P$ est \struck{les} torseur des
sections horizontales de $P^\infty_{S/T}(P)$, on voit que
\struck{cette condition} cette condition signifie que
\[
(28)\qquad \lambda(P) \subset \bigl(P^\infty_{S/T}(E)/E\bigr)^{\natural} ,
\]
où le \struck{\ill{}} signe $\natural$ désigne le sous-faisceau
des sections horizontales.
\marginal{\uncertain{petite} \uncertain{remarque} : \ill{} \uncertain{donc} $P^\infty_{S/T}(P)$ \uncertain{dans} un torseur stratifié quelc.}
\note{le signe noté $\natural$ est un petit éclair ou zigzag en exposant ; on ne lui connaît pas d'équivalent sûr.}

\struck{En} \emph{En résumé}, \struck{\ill{}} si $P$ est un $E$-torseur,
\begin{quote}
une stratification sur $E$ rel. à $T$, compatible avec
la stratification $\delta = \delta^\infty_E$ donnée sur $E$, équivaut à
un hom. de faisceaux (26), satisfaisant
(27) et (28) ; cette dernière condition permet d'ailleurs
d'exprimer en \uncertain{précisant} la donnée (26) \uncertain{par}
\[
(26')\qquad P \xrightarrow{\ \lambda\ } \bigl(P^\infty_{S/T}(E)/\delta E\bigr)^{\natural} .
\]
\end{quote}
\note{le « une stratification sur $E$ » de l'énoncé encadré : on attendrait « sur $P$ ».}

Cette description est fonctorielle en $E$, et
compatible aux changements de base.

\section{3) Relations entre stratifications et connexions intégrables sur torseurs affines}

\page{151}
Si $P$ est un torseur sous $E$, $E$ muni d'une
stratification $\delta^\infty_E = \delta$ rel. à $T$, \struck{\ill{}} toute stratification de $P$
\add{compatible à $\delta$} définit évidemment une connexion (compatible à $d^0_E$).
\struck{Je dis évidemment que l'hom. canonique (24), \ill{} compatible avec cette association}
Il convient encore de préciser quelle est l'hom.
$d^0_P : P \to E \otimes \Omega^1_{S/T}$ correspondant à la description (17).
Pour l'expliciter, notons que l'inclusion
\[
P^\infty_{S/T}(E)^+ \hookrightarrow P^\infty_{S/T}(E)
\]
fournit, par composition, un isomorphisme additif
\[
(29)\qquad P^\infty_{S/T}(E)^+ \simeq P^\infty_{S/T}(E)/\delta E ,
\]
l'hom. correspondant
\[
(29')\qquad P^\infty_{S/T}(E) \xrightarrow{\ p\ } P^\infty_{S/T}(E)/\delta E \xrightarrow{\ \sim\ } P^\infty_{S/T}(E)^+
\]
\note{une accolade sous la composée de (29').}
étant donné par
\[
\text{\struck{$d^+_E$}}\qquad \varphi \longmapsto \varphi - \delta_E(\varphi) .
\]
\note{ainsi sur la page ; on attendrait $\varphi - \delta_E(\varepsilon_E\varphi)$.}
Donc l'hom. $\lambda$ de (26) s'interprète comme un
\[
(30)\qquad \lambda^+ : P \longrightarrow P^\infty_{S/T}(E)^+ ,
\]
\struck{la condition $\lambda^+$} la condition (27) s'écrivant alors
\[
(31)\qquad \lambda^+(x + \xi) = \lambda^+(x) + \underbrace{\bigl[d^\infty_{S/T}(\xi) - \delta^\infty_E(\xi)\bigr]}_{d^+_{S/T}(\xi)} .
\]
\marginal{Définir les sections \emph{intégrables} de $P^\infty_{S/T}(E)^+$ comme celles correspondant aux sections
\struck{\ill{}} de $(P^\infty_{S/T}(E)/\delta E)^{\natural}$, et introduire le \uncertain{sous-faisceau} \ill{} de
$P^\infty_{S/T}(E)^+$. [Mais attention, $P^\infty_{S/T}(E)^+ \hookrightarrow P^\infty_{S/T}(E)$
\struck{\ill{}} n'est \emph{pas} compatible avec stratifications.]}

\page{152}
Réduisant à valeurs dans $P^1_{S/T}(E)^+ = E \otimes \Omega^1_{S/T}$, on
trouve un hom.
\[
(32)\qquad d^0_P : P \longrightarrow E \otimes \Omega^1_{S/T} , \qquad d^0_P(x) = \lambda^+(x) \bmod \ldots
\]
satisfaisant
\[
d^0_P(x + \xi) = d^0_P(x) + d^0_E(\xi) ,
\]
compte tenu de (23).
\note{le renvoi « (23) » est écrit sur un autre chiffre ; dans (32), l'indice du $d^0$ de droite est lu $E$ sur la page, on attendrait $P$.}

\begin{quote}
On vérifie que (32) n'est autre que l'hom.
qui définit la connexion de $P$ associée à sa stratification, \struck{\ill{}} que cette connexion est intégrable, enfin que
si $\xi$ classe \uncertain{dans} $H^1_{DR}(S,E)$ n'est autre que
l'image de celle de \struck{\ill{}} $(P,\lambda)$ dans $H^1_{\mathrm{strat}}(S,E_{\mathrm{strat}})$
par l'homomorphisme canonique (24).
\end{quote}
\marginal{faire la vérification}
\note{l'alinéa est marqué d'un double trait vertical dans la marge.}

\emph{Corollaire} Supposons que $T$ soit de car. nulle et
\begin{quote}
$\Omega^1_{S/T}$ plat [p.ex. $S$ formellement lisse sur $T$] ; on sait
que (24) est un \emph{isomorphisme} [ ]. \struck{\ill{}} Soit
$E$ quasi-coh. sur $S$ avec stratification $\delta = \delta^\infty_E$, $P$
un torseur sous $E$. Alors il y a correspondance
biunivoque entre les stratifications de $P$ compatibles
avec celle de $E$, et les connexions intégrables
sur $P$ compatibles avec celle de $E$. Donc
aussi, entre les applications (26) $P \xrightarrow{\lambda} P^\infty_{S/T}(E)/\delta E$ \struck{\ill{}}
[i.e. encore $\lambda^+ : P \to (P^\infty_{S/T}(E)^+)^{\natural}$] satisfaisant les
conditions (27), (28) [\struck{\ill{}} \add{équivalent}, (31)] \struck{\ill{}} \add{et à valeurs} \add{dans $P^\infty_{S/T}(E)^{+\natural}$}
\end{quote}
\marginal{NB Sous les hyp. faites, les stratifications sur $E$ correspondant bijectivement aux conn. intégrables sur $E$ \ldots}
\note{après « coh. sur $S$ », un « $P$ un torseur » biffé. La phrase du corollaire se poursuit p.~153.}

\page{153}
\begin{quote}
\struck{\uncertain{fait} $\lambda = p\lambda^+$}], et les hom. (17) $P \xrightarrow{d_P} E \otimes \Omega^1_{S/T}$
satisfaisant les conditions (18), en associant à $\lambda^+$
sa réduction dans $P^1_{S/T}(E)^+ \simeq E \otimes \Omega^1_{S/T}$.
\end{quote}

\emph{NB} Prenant \struck{$P$} triviale sur $E$, ceci \uncertain{prouve} que
l'hom. \struck{de réduction} de réduction $P^\infty_{S/T}(E)^+ \to P^1_{S/T}(E)^+$
induit un \struck{bijection} \emph{isomorphisme de faisceaux}
\[
P^\infty_{S/T}(E)^{+\natural} \xrightarrow{\ \sim\ } Z^1(E \otimes \Omega^\bullet_{S/T}) ,
\]
où \uncertain{prenant} \ill{} la deuxième \uncertain{avec} \uncertain{comme} stratification
sur $P^\infty_{S/T}(E)^+$ celle déduite de la strat. can. de
$P^\infty_{S/T}(E)/\delta E$ par transport de structure : l'aide de
(29). Ainsi, \emph{une} 1-\emph{forme fermée sur $S/T$}
\add{(\emph{coeff. dans} $E$)} \emph{définit un \uncertain{jet} d'ordre $\infty$} \add{$\lambda^+$} \emph{de $E$ sur $S/T$},
satisfaisant la condition (28). À la forme
exacte $d^0_E(\xi)$ est associé ainsi \struck{\ill{}}
\struck{\ill{}} $\lambda(\xi) = d^\infty_E(\xi)$ i.e. $\lambda^+(\xi) = d^\infty_E(\xi) - \delta^\infty_E(\xi) = d^+_E(\xi)$.
\note{dans la dernière formule, $\lambda$ est écrit sur un autre signe, et « $= d^\infty_E(\xi)$ » est lu sur un signe surchargé.}

\section{4) Équations de Picard-Fuchs}

Supposons encore qu'on ait $E$ sur $S$, avec
stratification $\delta^\infty_E = \delta$. Supposons donnée de plus
une suite exacte de modules quasi-coh.
\[
(33)\qquad 0 \longrightarrow E' \xrightarrow{\ u\ } E \xrightarrow{\ v\ } E'' \longrightarrow 0 .
\]
On ne suppose pas que $E'$ soit un sous-module
stratifié. \struck{\ill{}} Supposons donné un
\note{la phrase se poursuit en tête de la p.~155.}

\page{155}
torseur $P'$ sous $E'$, posons
\[
(34)\qquad P = P' \times^{E'} E .
\]
La donnée d'une stratification sur $P$ compatible avec
celle de $E$, grâce au \uncertain{n° 2}, équivaut à la donnée
d'un hom.
\[
(35)\qquad \lambda' : P' \longrightarrow \bigl(P^\infty_{S/T}(E)/\delta E\bigr)^{\natural}
\]
satisfaisant
\[
(36)\qquad \lambda'(x' + \xi') = \lambda'(x') + p\,d^\infty_E(u\xi') \qquad x' \in P'(\cdot),\ \xi' \in E'(\cdot) .
\]
\note{dans (36), « $p\,d^\infty_E$ » est récrit sur un premier essai surchargé.}
On peut aussi interpréter $\lambda'$ comme un
\[
(35')\qquad \lambda'^{+} : P' \longrightarrow \bigl(P^\infty_{S/T}(E)^+\bigr)^{\natural} , \qquad \text{(cf. transport de structure)}
\]
satisfaisant
\[
(36')\qquad \lambda'^{+}(x' + \xi') = \lambda'^{+}(x') + \underbrace{d^+_E(u\xi')}_{= d^\infty_E(u\xi') - \delta^\infty_E(u\xi')} .
\]

\struck{\ill{} la condition d'intégrabilité, signifiant} \add{de $p$-faisceaux}

\struck{Notons donc} Notons alors la suite exacte
\[
(37)\qquad \cdots \longrightarrow P^\infty_{S/T}(E') \xrightarrow{\ u\ } P^\infty_{S/T}(E) \xrightarrow{\ v\ } P^\infty_{S/T}(E'') \longrightarrow 0 ,
\]
compatible aux stratifications, \add{de $p$-faisceaux}.
Ceci posé, grâce à (35) \add{et (36)}, \uncertain{une} section canonique
\[
(38)\qquad \varphi = \varphi(\lambda') \in \Gamma\bigl(S, P^\infty_{S/T}(E'')/v\delta(E)\bigr)
\]
\[
\simeq \Gamma\bigl(S, P^\infty_{S/T}(E'')^+/v\delta u(E')\bigr) .
\]
\note{la place de « de $p$-faisceaux », écrit en marge droite avec un trait de renvoi, est incertaine ; dans (38) le dénominateur est surchargé et la lecture $v\delta u(E')$ n'est pas sûre.}
\struck{En outre}, $\varphi$ est \struck{\ill{}} même
\[
(38')\qquad \varphi \in \Gamma\Bigl(S, \bigl(P^\infty_{S/T}(E'')/v\delta(E)\bigr)^{\natural}\Bigr) \simeq \Gamma\bigl(S, P^\infty_{S/T}(E'')^{+\natural}/\delta u(E')\bigr) .
\]

\page{156}
On l'appelle la \struck{section} \emph{jet de Picard-Fuchs}
\emph{défini par} $P'$ \struck{\ill{}} et la \emph{stratification} \add{$\sigma$} \emph{donnée}.
\struck{\ill{}} $P = P' \times^{E'} E$. Ce n'est pas \uncertain{H} : \uncertain{faux} un
jet \add{intégrable} \struck{de} $E''$, mais \uncertain{seulement} \ill{} \uncertain{à} \ill{} \add{\uncertain{l'image} $v\delta^\infty_E u$}
c'est-à-dire un tel jet, défini \uncertain{modulo} \struck{\ill{}} \add{l'image} \add{d'une}
section de \struck{\ill{}} $E'$.

Soit $F$ un module \add{\uncertain{quasi-coh.}} sur \struck{$E$} $S$, et
\[
(39)\qquad D : E'' \longrightarrow F
\]
un opérateur différentiel \add{\uncertain{linéaire}} rel.~$/T$, définissant donc un
hom. \add{\uncertain{linéaire}} \struck{\ill{}} : \struck{\ill{}}
\[
(40)\qquad u_D : P^\infty_{S/T}(E'') \longrightarrow F .
\]
On dit que $u_D$ \struck{$D$} est un \emph{opérateur de Picard-Fuchs}, \emph{relativement} à la suite exacte
(33), \uncertain{à} l'application linéaire correspondante
(40) s'annule sur $v\delta(E)$, i.e. si $u_D$ se
factorise en un homomorphisme
\[
(41)\qquad \bar u_D : P^\infty_{S/T}(E'')/v\delta(E) \longrightarrow F .
\]
\struck{Soit alors} Si alors $(P',\sigma)$ est comme
ci-dessus, d'où un $\varphi$ (38), on peut
former $\bar u_D(\varphi) \in \Gamma(S,F)$ ; on pose
\[
(42)\qquad \text{\struck{$\ldots$}}\ \mu_D(P,\sigma) = \bar u_D\bigl(\varphi(P,\sigma)\bigr) \in \Gamma(S,F) ,
\]
\note{« On dit que $u_D$ » : $u_D$ est écrit sur un $D$, et le second « est » est lu tel quel. Dans (42), le premier membre biffé est illisible ; il écrit $(P,\sigma)$ là où l'on attendrait $(P',\sigma)$. La phrase se poursuit p.~157.}

\page{157}
on l'appellera la \emph{section de Manin de $F$} associée
à l'op. de Picard-Fuchs $D$ et aux données $(P',\sigma)$.
\emph{Équations de Pic.-Fuchs universelles} : à valeurs dans le pro-faisceau
$P^\infty_{S/T}(E'')/v\delta u(E')$ ; elle \uncertain{redonne} la section $\varphi$ dont il s'agit.
\note{deux traits horizontaux séparent ce qui précède de ce qui suit.}

Lorsque $E$ est seulement muni d'une connexion
intégrable \add{rel. à $T$}, \struck{\ill{}} (au lieu d'une stratification), on voit
que les connexions intégrables sur $P = P' \times^{E'} E$
correspondent aux hom.
\[
(43)\qquad P' \xrightarrow{\ d_{P'}\ } E \otimes \Omega^1_{S/T}
\]
\note{le but de $d_{P'}$ est écrit sur un premier symbole raturé.}
satisfaisant
\[
(44)\qquad
\begin{cases}
d_{P'}(x' + \xi') = d_{P'}(x') + d^0_E\bigl(u(\xi')\bigr) \\
d_{P'}(P') \subset Z^1(E \otimes \Omega^\bullet_{S/T}) , \ \text{i.e. } d^1_E \circ d_{P'} = 0 .
\end{cases}
\]

Dans ce cas, \uncertain{il y a} une variante de la
notion d'équation \add{de}/Picard-Fuchs, \add{\uncertain{pour} (30)} \struck{\ill{}} nous dirons les
opérateurs différentiels
\[
(45)\qquad E \otimes \Omega^1_{S/T} \xrightarrow{\ \Delta\ } F
\]
\note{l'indice de $\Delta$ est raturé.}
nuls sur \struck{l'image de $E' \otimes \Omega^1_{S/T}$}, \struck{i.e. les opérateurs}
\add{$d^0_E(u(E'))$}, ceux de la
[qui se factorisent \struck{\ill{}} \uncertain{pour} \uncertain{module} \struck{$E'' \otimes \Omega^1_{S/T} \to$}
forme $\Delta' \circ d^1_E$, \uncertain{où} $E \otimes \Omega^2_{S/T} \xrightarrow{\Delta'} F$ un
autre op. diff.]. Comme $d_{P'}$ (43) fournit
\add{\uncertain{localement}} une section de $E \otimes \Omega^1_{S/T}$ modulo une section
\note{la phrase se poursuit p.~158.}

\page{158}
de la forme $d^0_E(u(\xi'))$, $\xi' \in E'(\cdot)$, on voit que
$\Delta$ définit une section
\[
(46)\qquad \nu_\Delta(P', d_{P'}) \in \Gamma(S,F) .
\]
Ici la situation universelle est obtenue en
prenant pour $F$ le pro-faisceau
\[
F = P^\infty_{S/T}(E \otimes \Omega^1_{S/T})/I\bigl(P^\infty_{S/T}(E')\bigr) ,
\]
\note{la lecture « $I$ » devant $P^\infty_{S/T}(E')$ est incertaine.}
où $P^\infty_{S/T}(E')$ \uncertain{s'envoie} \uncertain{dans} $P^\infty_{S/T}(E \otimes \Omega_{S/T})$ provient de l'op.
diff. composé \add{$d^0_E u$ :} \add{(\uncertain{de degré} 1)} $E' \xrightarrow{u} E \xrightarrow{d^0_E} E \otimes \Omega^1_{S/T}$.

Revenons au cas où $E$ admet une stratification, \struck{\ill{}} \struck{et $P$ muni d'une stratification compatible avec}
\add{\uncertain{soit} $(P',\sigma)$} \struck{\ill{} pour \ill{}}
et prenons $d^0_E$ déduit d'icelle ; \struck{\ill{}} \struck{\ill{}}
\struck{(comme on \ill{}} \struck{\ill{} les sections de} \struck{\ill{} que \ill{}}. On compare la
section $\nu_\Delta$ avec $\mu_D$ de la façon suivante.
À tout $\Delta$ \struck{$\ldots$} est associé un opérateur de Picard-Fuchs,
\struck{\ill{}} pour (30), en prenant le composé
\[
(47)\qquad D_0 : E \xrightarrow{\ d^0_E\ } E \otimes \Omega^1_{S/T} \xrightarrow{\ \Delta\ } F ,
\]
et on prouve aisément que \add{en $D : E'' \to F$,} \struck{\ill{}} \uncertain{ce qui est}
possible puisque par hyp. sur $\Delta$, $\Delta \circ d^0_E$ est nulle
sur $E'$. Ceci posé, on trouve que
\[
(48)\qquad \nu_\Delta = \mu_D ,
\]
ce qui donne la réduction annoncée.
\marginal{faire la vérification}
\note{dans le renvoi « pour (30) », on attendrait plutôt (33) ; le numéro est lu tel quel.}

\page{159}
Dans le cas favorable envisagé à la fin du n° 3,
($T$ de car. nulle, $\Omega^1_{S/T}$ plat), on peut donc
exprimer les $\mu_D$ en termes des $\nu_\Delta$. Pour
ceci, il suffit de voir que, pour $S$ affine, chaque
$D$ provient bien d'un $\Delta$. Ceci résulte du
fait suivant (lui-même cas particulier d'un fait
bien plus général sur les op. diff.) : \struck{\ill{}}.

\emph{Lemme}. Sous les conditions précédentes ($S$ car.
\begin{quote}
nulle, $\Omega^1_{S/T}$ plat), \struck{les} op. diff. \struck{$E$} $E \to F$
qui sont tels que $P^\infty_{S/T}(E) \to F$ s'annule sur
$\delta(E)$, sont exactement ceux qui \struck{\ill{}} de la
forme $\Delta \circ d^0_E$, avec $\Delta : E \otimes \Omega_{S/T} \to F$ un
\struck{\ill{}} op. différentiel. L'indétermination est exactement par un
\struck{$\Delta' \circ d^1_E : E \otimes \Omega_{S/T}$} \add{$\Delta' \circ d^1_E$}, où $\Delta' : E \otimes \Omega^2_{S/T} \to F$
\uncertain{est un} op. diff.
\end{quote}
\struck{Ceci résulte du fait que \ill{} $E \to \Omega$}
\note{« $S$ car. nulle » est sur la page ; plus haut il écrit « $T$ de car. nulle ». Le lemme est marqué d'un trait vertical dans la marge. Dans la dernière phrase, $\Delta'$ est écrit sur un $\Delta$, et la lecture du premier membre biffé ($\Delta' \circ d^1_E$, $E \otimes \Omega_{S/T}$) est incertaine.}

\note{La page 160 est la couverture brune du dossier ; elle porte de sa main : « Modules stratifiés et \struck{Connexions intégrables et} équations de Picard-Fuchs ». Elle n'est pas transcrite comme page ; la suite des notes, si elle existe, est hors de ce lot.}

\end{document}
